\begin{summary}
\textbf{Question.} Do paragliders move differently when they fly in a group? If they do, is it the group, or
who flies in groups, where and when they fly, or the moment they choose to fly together?

\textbf{Answer.} Over a quarter of an hour the progress in a group is straighter: $\Delta H_\mathrm{loc}(1000)$ =
\DHlocThousandRaw{} raw, \DHlocThousandFlight{} within the same flight, \DHlocThousandGroupSdh{} at the same site-day and
hour, and only \DHlocThousandSel{} from who flies in groups. It is not the zone (reweighting removes at most
\ZoneRedLocZoneHour\% of it) and not the moment (on the same flight it starts with the group, \PPEntryDHloc, and stops
with it, \PPExitDHloc). Within one glide or climb (100~s) nothing changes.

\textbf{Also in a group,} once stratified: more time climbing, more searches that end in a climb, a higher climb rate,
faster and flatter glides, a higher cross-country speed.

\textbf{Inside the group.} Most of the effect is there with two pilots; at the same place and hour only 7+ stands apart.
The pilot ahead flies a little faster but finds no more thermals.

\textbf{Limits.} Strata and windows fix pilot, day, site, hour, zone and the drift of the flight, not what brings pilots
together; a contact is shared air within sight; phases exist on 1~Hz flights only. Data: \tauStar-min contacts of the
note \emph{Group flights in free flight}, 154\,598 paraglider flights (21.7\% in a group at some point).
\end{summary}
